\scanpage{40}
On conclut du premier groupe
\[
\varphi_*\Lambda_Y\varphi^*
 =\underbrace{\varphi_*\varphi^*\Lambda_X}_{\substack{L_X\Lambda_X\\ \shortparallel\\ \mathrm{id}_X-\sum_{0\leq i\leq\unsure{n|2n}}p_i^X}}
 \left(\mathrm{id}_X-\sum_{n+1\leq j\leq2n}p_X^j\right)
 =\underbrace{\mathrm{id}_X-\sum_{0\leq j\leq2n}p_X^j}_{q}.
\]
$[$Plus gén.
\[
\varphi_*\Lambda_Y^m\varphi^*
=\deleted{\unsure{\sum|\Sigma}}\left(\mathrm{id}-\sum_{0\leq i\leq n}p_X^i\right)\Lambda_X^{m-1}
 \left(\mathrm{id}-\sum_{n+1\leq j\leq2n}p_X^j\right)=q\Lambda_X^{m-1}q
\]
où $q=\mathrm{id}-\sum_{0\leq i\leq2n}p_X^i$

Aussi :
\[
\varphi^*\Lambda_X\varphi_*=\mathrm{id}_Y-p_E
\]
\inserted{projection sur partie évanescente $E^{n-1}$}

d'où, en multipliant par $\Lambda_Y$ à gauche, comme $\Lambda_Yp_E=0$
\[
\Lambda_Y=\underline{\Lambda_Y\varphi^*}\Lambda_X\varphi_*=
 \underline{\varphi^*\Lambda_X\Bigl(\mathrm{id}_X-\sum_{n+1\leq j\leq2n}}
 \overbrace{\underline{p_X^j\Bigr)}\,\Lambda_X}^{\substack{0\\\shortparallel}}\varphi_*\qquad\text{donc}
\]
\[
\boxed{\Lambda_Y=\varphi^*\Lambda_X^2\varphi_*}
\qquad\text{Donc}\quad C(X)\quad\text{\deleted{\unsure{équivaut|implique}}\inserted{implique}}\quad C(Y).
\]

Soit $p_X^{i,j}$ \unsure{la|les} projection sur $L_X^{\unsure{j|t}}P^i(X)$ \quad $(i+j\leq n,\ i,j\geq0)$.

Alors
\margin{$\unsure{\lambda|L|\Lambda}$?}
\[
\varphi_*p_Y^{i,j}\inserted{\unsure{\lambda|\Lambda}}\varphi^*=p_X^{i,j+1}
 \qquad\text{pour}\qquad \unsure{0|i}+j\unsure{<|\leq}n-1,\quad i,j\geq0\quad\deleted{\unsure{\text{algébr}}}
\]
donc la condition $D(Y)$ implique que les $p_X^{ij}$ sont algébriques pour $\deleted{\unsure{j}}\,i+j<n$, $i\geq0$ et \uline{$j\geq1$}, i.e. les seuls \unsure{évt. non alg.|int. en susp.|rest. en susp.} sont les \uline{$p_X^i$ pour $0\leq i\leq\unsure{2n|n|\deleted{2n}}$}. Mais
% archive sheet number: 39

